%
% vox-chunked-transfer.tex
%
% Z Specification for the vox daemon's chunked `get' (fetch) transport.
%
% The daemon delivers a stored file of arbitrary total size in bounded
% per-frame memory: fetch_begin, then ordered `chunk' frames, then fetch_end
% (or an error terminal). This model is deliberately kept out of
% audio-programs.tex: it touches no Program, Radio, or Catalog field. It is
% transport, on the order of the record-store containment model
% (vox-dvri-record-store.tex), and pins the four integrity properties a
% mid-stream fault must not violate: ordered contiguous chunks, a reassembled
% total equal to the declared byte count, a containment check taken once before
% the first byte, and an atomic terminal --- the complete file lands or nothing
% does.
%

\documentclass[a4paper,10pt,fleqn]{article}
\usepackage[margin=1in]{geometry}
\usepackage{fuzz}
\usepackage[colorlinks=true,linkcolor=blue,citecolor=blue,urlcolor=blue]{hyperref}

\begin{document}

\title{Vox Chunked Transfer: A Z Specification}
\author{Formal model of the daemon \texttt{get}/\texttt{fetch} transport}
\date{July 2026}
\hypersetup{
  pdftitle={Vox Chunked Transfer: A Z Specification},
  pdfauthor={Punt Labs},
  pdfsubject={Z Specification},
  pdfcreator={fuzz/probcli}
}
\maketitle

\section{Introduction}

The daemon's \texttt{get} verb retrieves a store file --- a recording, or one
part of a music album --- over the single client socket. The old transport put
the whole file in one base64 frame and capped the total at
$700\,000$ bytes; the redesign (\texttt{docs/rec-music-cli.md},
\S``Chunked transfer'') replaces it with a bounded stream:

\begin{quote}
\texttt{fetch\_begin} $\rightarrow$ \texttt{chunk}$^{*}$ $\rightarrow$
(\texttt{fetch\_end} $\mid$ \texttt{error})
\end{quote}

Each \texttt{chunk} carries at most \texttt{FETCH\_CHUNK\_BYTES}; the file may be
arbitrarily large. The client reassembles the chunks in \texttt{seq} order,
checks the running total against the declared \texttt{bytes}, and writes the
result atomically, so a mid-stream fault never leaves a partial file. This
specification models exactly that discipline, and states the four integrity
properties as invariants and preconditions rather than as prose the code is
merely asked to honour.

The properties, and where each lives in the model:

\begin{itemize}
  \item \textbf{Ordered, contiguous chunks.} \texttt{Chunk} requires
    $seq? = next$ and advances $next$ by one, so the delivered sequence numbers
    are exactly $0, 1, \dots$ with no gap, reorder, or repeat.
  \item \textbf{Reassembled total equals declared bytes.} \texttt{End} is
    enabled only when $got = declaredBytes$ (and all declared chunks have
    arrived); the running total $got$ never exceeds $declaredBytes$ by the state
    invariant.
  \item \textbf{Containment checked once, before the first byte.}
    \texttt{Begin} is the sole transition into the streaming state and it is the
    only place the reference is checked; an unsafe reference is refused by
    \texttt{RejectBeforeBegin} from the idle state, so no byte is ever sent for
    an uncontained reference.
  \item \textbf{Atomic terminal.} A file is marked \emph{landed} only in the
    \texttt{delivered} phase, which requires the complete byte count; a
    mid-stream fault takes \texttt{Abort} to the \texttt{failed} phase with
    nothing landed. No reachable state carries a landed partial file.
\end{itemize}

\section{Basic Types}

A \texttt{REF} is a client-supplied reference to a store member: a recording id,
or a music part name resolved inside its album directory. The model treats it as
opaque; the containment rule of \texttt{vox-dvri-record-store.tex} decides which
references are safe, abstracted here as a set. A \texttt{REASON} is an opaque
non-empty diagnostic string, as in the audio-programs model.

\begin{zed}
  [REF, REASON]
\end{zed}

A Z boolean, spelled to avoid the ProB/B keyword clash with \texttt{BOOL}. The
\texttt{landed} flag records whether the complete file has been placed on disk.

\begin{zed}
  ZBOOL ::= ztrue | zfalse
\end{zed}

A transfer is in one of four phases. \texttt{idle} is before any request;
\texttt{active} is after \texttt{fetch\_begin}, receiving chunks;
\texttt{delivered} is the success terminal, reached by \texttt{fetch\_end};
\texttt{failed} is the error terminal, reached by an \texttt{error} frame. The
two terminals are distinct and absorbing: one fetch is one lifecycle on one
socket.

\begin{zed}
  Phase ::= idle | active | delivered | failed
\end{zed}

\section{Global Constant}

\texttt{chunkLimit} is the per-chunk byte bound (\texttt{FETCH\_CHUNK\_BYTES} in
the design; e.g.\ 256\,KiB). Its exact value is an implementation tuning knob;
what the model fixes is that it is a positive constant and that every chunk
respects it. There is deliberately \emph{no} total-file bound: the ceiling the
old transport imposed is gone.

\begin{axdef}
  chunkLimit : \nat_1
\end{axdef}

\section{Containment}

Which references resolve safely under the daemon-owned root is decided by the
containment rule modelled in full in \texttt{vox-dvri-record-store.tex} (absolute,
traversing, separator-bearing, empty, or non-printable names are rejected before
any filesystem touch). Here that rule is abstracted as the set of safe
references; the point of this model is \emph{when} the check happens, not how it
is computed.

\begin{axdef}
  contained : \power REF
\end{axdef}

\section{State}
\label{sec:state}

The \texttt{Transfer} state carries the phase, the declared totals from
\texttt{fetch\_begin} (\texttt{declaredBytes}, the \texttt{bytes} field, and
\texttt{declaredChunks}, the \texttt{chunks} field), the next expected sequence
number \texttt{next} (which equals the number of chunks received so far), the
running reassembled total \texttt{got}, and the \texttt{landed} flag.

The predicate carries every integrity invariant. The running total never
overshoots the declaration and never more chunks arrive than declared; a file is
\emph{landed} exactly in the \texttt{delivered} phase, and that phase is reached
only with the full byte count and every chunk present; the \texttt{failed} phase
never lands anything; and \texttt{idle} is the pristine zero state.

\begin{schema}{Transfer}
  phase : Phase \\
  declaredBytes : \nat \\
  declaredChunks : \nat \\
  next : \nat \\
  got : \nat \\
  landed : ZBOOL
  \where
  got \leq declaredBytes \\
  next \leq declaredChunks \\
  (landed = ztrue \iff phase = delivered) \\
  (phase = delivered \implies (got = declaredBytes \land next = declaredChunks)) \\
  (phase = failed \implies landed = zfalse) \\
  (phase = idle \implies \\
  \quad~ (declaredBytes = 0 \land declaredChunks = 0 \\
  \quad~ \land next = 0 \land got = 0 \land landed = zfalse))
\end{schema}

The invariant $landed = ztrue \iff phase = delivered$ together with
$phase = delivered \implies got = declaredBytes$ is the heart of the atomic
terminal: a landed file is always the complete file. No state schema admits a
\texttt{landed} partial, so the property holds by construction, not by a runtime
check the client must remember.

\section{Initialisation}

A fresh transfer is idle with all counters zero.

\begin{schema}{InitTransfer}
  Transfer'
  \where
  phase' = idle \\
  declaredBytes' = 0 \\
  declaredChunks' = 0 \\
  next' = 0 \\
  got' = 0 \\
  landed' = zfalse
\end{schema}

\section{Operations}

\subsection{Begin --- the \texttt{fetch\_begin} frame}

From \texttt{idle}, the daemon resolves and containment-checks the reference
\emph{once} and announces the totals. The declared chunk count must be able to
carry the declared bytes ($declaredBytes \leq declaredChunks * chunkLimit$) and
must not exceed them ($declaredChunks \leq declaredBytes$, so each declared chunk
holds at least one byte); together these force $declaredBytes = 0 \iff
declaredChunks = 0$, the empty-file case. The reference is checked here and only
here: \texttt{Begin} is the sole gateway into \texttt{active}, so a byte is sent
only after containment has passed.

\begin{schema}{Begin}
  \Delta Transfer \\
  ref? : REF \\
  bytes? : \nat \\
  chunks? : \nat
  \where
  phase = idle \\
  ref? \in contained \\
  bytes? \leq chunks? * chunkLimit \\
  chunks? \leq bytes? \\
  phase' = active \\
  declaredBytes' = bytes? \\
  declaredChunks' = chunks? \\
  next' = 0 \\
  got' = 0 \\
  landed' = zfalse
\end{schema}

\subsection{RejectBeforeBegin --- the pre-stream error frame}

An unsafe reference is refused from \texttt{idle} with an error frame emitted
\emph{before} any \texttt{fetch\_begin}. The state does not move: no totals are
announced, no byte crosses the wire, nothing is landed. This is the model image
of ``a hostile ref is audit-logged and rejected before the first byte''.

\begin{schema}{RejectBeforeBegin}
  \Xi Transfer \\
  ref? : REF
  \where
  phase = idle \\
  ref? \notin contained
\end{schema}

\subsection{Chunk --- a \texttt{chunk} data frame}

While \texttt{active} and with chunks still owed, the next data frame must carry
the next sequence number ($seq? = next$) --- this is the ordered, contiguous,
no-repeat guarantee --- and at most \texttt{chunkLimit} bytes, at least one. The
running total grows by the frame length and may not overshoot the declaration.

\begin{schema}{Chunk}
  \Delta Transfer \\
  seq? : \nat \\
  len? : \nat
  \where
  phase = active \\
  next < declaredChunks \\
  seq? = next \\
  1 \leq len? \land len? \leq chunkLimit \\
  got + len? \leq declaredBytes \\
  phase' = active \\
  declaredBytes' = declaredBytes \\
  declaredChunks' = declaredChunks \\
  next' = next + 1 \\
  got' = got + len? \\
  landed' = zfalse
\end{schema}

Because \texttt{next} starts at $0$ (from \texttt{Begin}) and each \texttt{Chunk}
advances it by one while demanding $seq? = next$, the accepted sequence numbers
are precisely $0, 1, \dots, declaredChunks - 1$: strictly increasing, contiguous,
and each used once. A frame out of order, duplicated, or skipping a number has
$seq? \neq next$ and is not a legal \texttt{Chunk}.

\subsection{End --- the \texttt{fetch\_end} frame, atomic land}

The success terminal fires only when every declared chunk has arrived
($next = declaredChunks$) and the reassembled total equals the declared byte
count ($got = declaredBytes$). Only then is the file \emph{landed}. This is the
atomic write: the complete file appears, in one step, or --- via \texttt{Abort}
--- not at all.

\begin{schema}{End}
  \Delta Transfer
  \where
  phase = active \\
  next = declaredChunks \\
  got = declaredBytes \\
  phase' = delivered \\
  declaredBytes' = declaredBytes \\
  declaredChunks' = declaredChunks \\
  next' = next \\
  got' = got \\
  landed' = ztrue
\end{schema}

The precondition $got = declaredBytes \land next = declaredChunks$ is exactly the
client-side reassembly check ``the running total against \texttt{bytes}'', made a
guard on the transition that lands the file rather than an afterthought. A short
final chunk, or a dropped chunk, leaves $got < declaredBytes$ or $next <
declaredChunks$, so \texttt{End} is disabled and the only terminal available is
\texttt{Abort}.

\subsection{Abort --- the mid-stream \texttt{error} terminal}

A read fault, a short read, or any mid-stream failure ends the exchange with an
\texttt{error} frame instead of \texttt{fetch\_end}. The phase moves to
\texttt{failed} and nothing is landed: the client discards whatever partial it
had accumulated. \texttt{Abort} is enabled from \texttt{active} whatever the
running total, so it models a fault at any point in the stream.

\begin{schema}{Abort}
  \Delta Transfer \\
  reason? : REASON
  \where
  phase = active \\
  phase' = failed \\
  declaredBytes' = declaredBytes \\
  declaredChunks' = declaredChunks \\
  next' = next \\
  got' = got \\
  landed' = zfalse
\end{schema}

Because $landed'= zfalse$ and the \texttt{failed} invariant keeps it so, no file
is placed for an aborted transfer even when $got$ was one chunk short of
$declaredBytes$. The partial is never on disk; ``the complete file lands or
nothing does'' is a fact of the two terminals, not of client discipline.

\section{Key Properties}

The following hold by construction (state invariants) or as operation
preconditions; they are the intended proof obligations for animation, and each
corresponds to a named property in \texttt{docs/rec-music-cli.md}
\S``Chunked transfer''.

\begin{itemize}
  \item \textbf{Ordered contiguous delivery}: every \texttt{Chunk} demands
    $seq? = next$ and sets $next' = next + 1$ from an initial $next = 0$, so the
    accepted sequence numbers are $0 \dots declaredChunks - 1$ with no gap,
    reorder, or repeat.
  \item \textbf{Bounded per-chunk memory}: every \texttt{Chunk} obeys
    $len? \leq chunkLimit$, and there is no total-file bound anywhere in the
    model --- a file of any size transfers.
  \item \textbf{Reassembled total equals declared bytes}: $got \leq
    declaredBytes$ always, and \texttt{End} requires $got = declaredBytes$ with
    $next = declaredChunks$, so a delivered transfer reassembled every declared
    byte in every declared chunk.
  \item \textbf{Containment once, before the first byte}: \texttt{Begin} is the
    only transition into \texttt{active} and the only reader of
    \texttt{contained}; an uncontained reference takes \texttt{RejectBeforeBegin}
    and never reaches \texttt{active}, so no chunk is ever emitted for it.
  \item \textbf{Atomic terminal}: $landed = ztrue \iff phase = delivered$, and
    $phase = delivered \implies got = declaredBytes$, so a landed file is always
    complete; \texttt{Abort} reaches \texttt{failed} with $landed' = zfalse$, so
    a mid-stream fault lands nothing. There is no reachable state with a landed
    partial file.
  \item \textbf{Distinct terminals}: \texttt{delivered} and \texttt{failed} are
    reached only by \texttt{End} and \texttt{Abort} respectively, and neither has
    an outgoing operation --- one request is one lifecycle, terminating in
    exactly one of success or error.
\end{itemize}

\end{document}
